\scanpage{78}\mspage{7\deleted{\unsure{6}}7}
posé, si $X'_K\mathrel{>\mkern-14mu<}Y'_K$, alors $X'\mathrel{>\mkern-14mu<}Y'$. En d'autres termes, si $\forall\, X''\in\widetilde{X'}$ \unsure{tels|tel} que $X''\ll K$ ($\Leftrightarrow X''\ll Y$), et tout $Y''\in\widetilde{Y'}$ \unsure{tels|tel} que $Y''\ll K$ ($\Longleftrightarrow Y''\ll X$), on a $X''\mathrel{>\mkern-14mu<}Y''$, alors $X'\mathrel{>\mkern-14mu<}Y'$.

Ceci implique At $11'$, \inserted{\deleted{\unsure{At 11}}} At 10, et la prop. de la page 73.
% slanted
\margin{\unsure{donc on peut} faire \unsure{sauter} At 10,}
Donc cela implique At 11 pour \deleted{familles \uline{finies}} $\widetilde F$ \uline{fini}, et le cor. 2 page 69 \unsure{pour} \unsure{ens|ens.} d'indices \uline{fini}.

\deleted{\unsure{Nous avons} un énoncé qui implique At 11} Les axiomes At 9 (critère de disjonction par les lieux) \inserted{\unsure{ses variantes} $9',9''$ (critères de compatibilité \inserted{par les lieux})} At 10 \unsure{.|'} (\inserted{premier} critère de disjonction des lieux), At 11 (critère de recollement des raffinements), \inserted{sa variante un peu affaiblie At $11'$,} \unsure{A12|At 12} (critère de compatibilité par intersection) sont de la nature de critères d'existence de figures (figures sommes, \deleted{\unsure{et}} en les quatre cas) intimement reliés entre eux, par des implications diverses\unsure{,|;} qu'on va résumer \unsure{des|dans} diagrammes. Mais je vais \inserted{d'abord} donner une formulation commune \deleted{\unsure{à}} \deleted{\unsure{qui|que}} à At $9'$ et At $9''$, par la \unsure{notion|création} de $\Sigma$-Atelier, où $\Sigma$ est une sous-cat. pleine de la cat. des ens. ordonnés. Un atelier est un $\Sigma$-atelier si
\[
\text{(i)}\quad \forall\,\deleted{\unsure{F}}\unsure{X|F}\in\deleted{\unsure{\mathcal F|\mathscr F}}M\qquad \widetilde{\deleted{F}X}\in\Sigma
\]
(ii) $\mathcal F$ satisfait l'axiome qui suit :
